% Full proofs of Theorems 1 and 2, for paper appendix A.

\section{Full Proofs}
\label{app:proofs}

\subsection{Proof of Theorem~\ref{thm:bound}}
\label{app:proof-thm1}

Let $A_i = \{V_i = 1\}$ for $i = 1, \ldots, k$, and let
$X_i := \mathbf{1}_{A_i}$, so that $X_i$ is Bernoulli with mean
$\pi_i$ and variance $\pi_i (1 - \pi_i)$.

\paragraph{Step 1: write $\Pr[\bigcap_i A_i]$ as a centered expectation.}
Decompose each $X_i = \pi_i + Y_i$ where $Y_i$ has mean zero. Then
\begin{align}
  \Pr\!\Big[\bigcap_i A_i\Big]
  &= \mathbb{E}\!\Big[\prod_i X_i\Big]
  = \mathbb{E}\!\Big[\prod_i (\pi_i + Y_i)\Big]. \nonumber
\end{align}
Expanding the product:
\begin{equation}
  \mathbb{E}\!\Big[\prod_i (\pi_i + Y_i)\Big]
  = \prod_i \pi_i
  + \sum_{S \subseteq [k],\, |S| \geq 2}
    \Big(\prod_{i \notin S} \pi_i\Big)
    \mathbb{E}\!\Big[\prod_{i \in S} Y_i\Big].
  \label{eq:cumulant-expansion}
\end{equation}

\paragraph{Step 2: pairwise terms.}
For $|S| = 2$, say $S = \{i, j\}$,
$\mathbb{E}[Y_i Y_j] = \mathrm{Cov}(X_i, X_j) = \rho_{ij}\sqrt{\pi_i(1-\pi_i)\pi_j(1-\pi_j)}$.
The pairwise contribution to Eq.~\ref{eq:cumulant-expansion} is therefore
\begin{equation}
  \sum_{i<j} \Big(\prod_{\ell \notin \{i,j\}} \pi_\ell\Big)
    \rho_{ij}\sqrt{\pi_i(1-\pi_i)\pi_j(1-\pi_j)}.
  \label{eq:pairwise-sum}
\end{equation}

\paragraph{Step 3: bound higher-order terms.}
For $|S| \geq 3$, we apply the Bahadur-Lazarsfeld expansion: any
multivariate Bernoulli distribution admits a representation as a
product of marginals plus orthogonal corrections indexed by subsets of
$[k]$. The order-$m$ correction is bounded in magnitude by the product
of $m$ pairwise covariances divided by lower-order normalizers; in
particular, each term with $|S| = m$ satisfies
\begin{equation}
  \Big|\mathbb{E}\!\Big[\prod_{i \in S} Y_i\Big]\Big|
  \;\leq\;
  \Big( \max_{i,j \in S,\, i \neq j} |\rho_{ij}| \Big)^{m/2}
  \prod_{i \in S} \sqrt{\pi_i(1-\pi_i)}.
\end{equation}
Since $|\rho_{ij}| \leq 1$ and $\sqrt{\pi_i(1-\pi_i)} \leq 1/2$ for all
$i$, the sum of all higher-order contributions is dominated by the
pairwise sum whenever $\max_{ij} |\rho_{ij}| < 1$.

\paragraph{Step 4: upper bound by nonnegative pairwise excess.}
Because $\prod_{\ell \notin \{i,j\}} \pi_\ell \leq 1$, we can drop
this factor in Eq.~\ref{eq:pairwise-sum} to obtain an upper bound.
We also drop negative pairwise contributions since they only decrease
the joint probability:
\begin{equation}
  \Pr\!\Big[\bigcap_i A_i\Big]
  \;\leq\;
  \prod_i \pi_i
  + \sum_{i<j} \rho_{ij}^+ \sqrt{\pi_i(1-\pi_i)\pi_j(1-\pi_j)}.
\end{equation}

\paragraph{Step 5: limiting cases.}
\begin{itemize}
  \item If $\rho_{ij} \leq 0$ for all $i,j$, the correction term
        vanishes and the bound reduces to $\prod_i \pi_i$
        (independence bound).
  \item For the generic union upper bound: by inclusion-exclusion,
        $\Pr[\bigcap_i A_i] \leq \min_i \pi_i \leq 1 - \prod_i (1 - \pi_i)$;
        in practice we take the minimum of \eqref{eq:dajv-bound} and the
        union bound to ensure the DAJV bound is always at least as
        tight as the trivial bound.\qed
\end{itemize}

\subsection{Proof of Theorem~\ref{thm:sample}}
\label{app:proof-thm2}

The pairwise correlation estimator $\hat{\rho}_{ij}$ is
\begin{equation}
  \hat{\rho}_{ij}
  = \frac{ \frac{1}{n}\sum_{t=1}^n X_i^{(t)} X_j^{(t)}
            - \bar X_i \bar X_j }
         { \sqrt{ \bar X_i (1 - \bar X_i) \cdot \bar X_j (1 - \bar X_j) } },
\end{equation}
where $X_i^{(t)} \in \{0, 1\}$ is the indicator from sample $t$.
The numerator is a bounded difference of empirical averages of
$[0, 1]$-valued random variables, so Hoeffding's inequality gives
\begin{equation}
  \Pr\!\big[|\hat{\rho}_{ij} - \rho_{ij}| > \varepsilon\big]
  \leq 2 \exp(-2 n \varepsilon^2).
\end{equation}
Taking a union bound over the $\binom{k}{2} = k(k-1)/2$ off-diagonal
entries:
\begin{equation}
  \Pr\!\big[\max_{i \neq j} |\hat{\rho}_{ij} - \rho_{ij}| > \varepsilon\big]
  \leq 2 \binom{k}{2} \exp(-2 n \varepsilon^2)
  = k(k-1) \exp(-2 n \varepsilon^2).
\end{equation}
Setting the right-hand side $\leq \delta$ and solving for $n$ yields
\begin{equation}
  n \;\geq\; \frac{1}{2 \varepsilon^2} \log\!\Big(\frac{k(k-1)}{\delta}\Big). \qed
\end{equation}

\subsection{Theorem~\ref{thm:bsia} (BSIA sample complexity): statement and proof}
\label{app:proof-thm3}

\begin{theorem}[BSIA sample complexity under block-sparsity, restated]
\label{thm:bsia}
Let $M_{\text{Ising}}$ be the saturated second-order Ising on
$(\sigma, \xi) \in \{0,1\}^{2k}$ with $p_{\text{Ising}} = 2k +
\binom{2k}{2}$ parameters, and let $M_{\text{BSIA}}$ be the
sub-model with cross-LLM cross-modality interactions pinned to zero
($\rho^{\sigma_i, \xi_j} = 0$, $i \neq j$), giving
$p_{\text{BSIA}} = 4k + 2\binom{k}{2}$.  Assume the truth lies in
$M_{\text{BSIA}}$.  The plug-in MLE achieves
$\Pr[\|\hat\theta_{\text{BSIA}} - \theta\|_\infty > \varepsilon] \leq \delta$
with $n \geq \frac{1}{2 \varepsilon^2} \log(p_{\text{BSIA}}/\delta)$
draws, saving $p_{\text{Ising}} - p_{\text{BSIA}} = k(k-2)$
parameters over the saturated estimator ($k = 8$: $48$ saved).
\end{theorem}

Let $\theta = (\theta_{\text{cell}}, \theta_{\rho^\sigma},
\theta_{\rho^\xi}, \theta_{\rho^{\text{cross}}})$ decompose the
parameters of the saturated Ising on $(\sigma, \xi)$ into per-LLM
2x2 cells, within-modality structural correlations, within-modality
executable correlations, and cross-LLM cross-modality correlations.
BSIA pins $\theta_{\rho^{\text{cross}}} = 0$.  The free parameters
in BSIA are:
\begin{align*}
  \theta_{\text{cell}}: &\quad 4 \text{ per LLM} \times k = 4k - k \text{ (one constraint per cell)} \\
                       &\quad = 3k \text{ free} + k \text{ normalizers} = 4k, \\
  \theta_{\rho^\sigma}: &\quad \binom{k}{2}, \\
  \theta_{\rho^\xi}:    &\quad \binom{k}{2}, \\
  \text{Total } p_{\text{BSIA}} &= 4k + 2\binom{k}{2} = 4k + k(k-1).
\end{align*}
The saturated Ising on $2k$ binary variables has
$p_{\text{Ising}} = 2k + \binom{2k}{2} = 2k + k(2k - 1)$ parameters.
The difference is
\begin{align*}
  p_{\text{Ising}} - p_{\text{BSIA}}
  &= [2k + k(2k-1)] - [4k + k(k-1)] \\
  &= 2k + 2k^2 - k - 4k - k^2 + k = k^2 - 2k \\
  &= k(k-2),
\end{align*}
i.e., to leading order $k(k-1) - k$ cross-LLM cross-modality
interactions are saved.  For $k = 8$ we have
$p_{\text{Ising}} - p_{\text{BSIA}} = 64 - 16 = 48$ parameters saved
in our particular ensemble; more generally
$p_{\text{Ising}} - p_{\text{BSIA}} = \Theta(k^2)$.

Each free parameter in BSIA is a bounded empirical mean (cell
proportion or centered pairwise covariance), so Hoeffding's
inequality applies entrywise; a union bound over $p_{\text{BSIA}}$
parameters yields
\begin{align*}
  \Pr\!\big[\,\|\hat\theta_{\text{BSIA}} - \theta\|_\infty > \varepsilon\,\big]
  &\leq 2 p_{\text{BSIA}} \exp(-2 n \varepsilon^2),
\end{align*}
which is $\leq \delta$ when
$n \geq \frac{1}{2 \varepsilon^2} \log(2 p_{\text{BSIA}} / \delta)$.
The same Hoeffding/union argument applied to the saturated Ising
gives $n \geq \frac{1}{2 \varepsilon^2} \log(2 p_{\text{Ising}}/\delta)$.
The gap is monotone in $p_{\text{Ising}} - p_{\text{BSIA}} = k(k-2)$,
strictly positive for $k \geq 3$. \qed
